\chapter{Primitives}

\section{Construction for flagging}\label{appendix:flags}

Let $\TF$ be a transformer, we define $\TF'$ that behaves like $\TF$ but recognize when a vector comes from embedding the token $\sigma$.
$\TF'$ has embedding dimension $d' = d + 1$ where $d$ is the embedding dimension of $\TF$.

Essentially, the components of $\TF'$ are the same as $\TF$.
The main change is the token embedding function $\TE' : \TF' : \Sigma \to \Floating^{d+1}$ which is defined as

\begin{equation*}
    (\theta'_{TE}(\tau))_j = \begin{cases}
        (\theta_{TE}(\tau))_j  & \text{ if } j \neq d+1 \\
        f_{\ON}                & \text{ if } j = d+1  \text{ and } \tau = \sigma \\
        f_{\OFF}               & \text{ if } j = d+1  \text{ and } \tau \neq \sigma
    \end{cases}
\end{equation*}

\noindent where $f_{\ON}$ is the value that marks the vector as flagged and $f_{\OFF}$ as unflagged.
Those values can be defined chosen arbitrarily for each flag.
In this case, the coordinate inserted for the flag is the last one, but it can be added anywhere.

The other components of $\TF'$ ignore the new $d+1$-th coordinate.
The position encoding of $\TF'$ is equal to the position encoding of $\TF$, except it puts a $0$ in the $d+1$ coordinate (not present in $\TF$) independently of the position.
The matrices of the $\ATTN$, $\FF$ and $\OUTPUT$ layers of $\TF'$ are the ones of $\TF$ extended with $0$ in the $d+1$ rows and columns to match the change in embedding dimension.

$\TF'$ behaves exactly as $\TF$ except it has one more coordinate in the embedding.
Said coordinate represents the flag and takes only the values $v_{\ON}$ or $v_{\OFF}$ in every vector.


\section{Ignore vectors in Attention layer}\label{appendix:ignore-vectors-attn}

Let $M$ be a set of indices.
We construct an attention layer of $\TF'$ such that $h_l$ ignore all the vectors $h_i$ such that $i\in M$.
As mentioned in the main text, we require the attention scores (previous to the $\softmax$) to satisfy the following:

\[
    \langle q_i', k_j' \rangle = \begin{cases}
        \minrep                  & \text{ if } i = l \land j \in M \\
        \langle q_i, k_j \rangle & \text{ if } i \neq l \lor j \notin M
    \end{cases}
\]

We achieved this by saturating the finite representation since $\langle q'_l, k'_i \rangle = \sum_{x=1}^{d+3} (q'_l)_x (k'_i)_x$.

\medskip

We force $(k'_i)_{d+2}$ and $(k'_i)_{d+3}$ to be $\minrep$ when $i \in M$, and $(q'_i)_{d+2}$ and $(q'_i)_{d+3}$ to be $1$ when $i=l$ and $0$ otherwise.
We also make $(k'_i)_{d+1}$ and $(q'_i)_{d+1}$ to be $0$ for all $i$ so the third to last term of the inner product is always $0$.
With these values, the only attention scores modified are the ones of the $l$th vector since every other will have $0$ in the last three coordinates of the query thus making the last three terms of the inner product $0$ and therefore not affecting the score.

This can be achieved with the following flags:
%
\begin{align*}
    & (h_i)_{d+1} = \begin{cases}
        1 & \text{ if } i = l \\
        0 & \text{ otherwise }
    \end{cases}
    & (h_i)_{d+2} = (h_i)_{d+3} = \begin{cases}
        1 & \text{ if } i \in M \\
        0 & \text{ otherwise }
    \end{cases}
\end{align*}
%
Now the query and key matrices are extended as:
%
\begin{align*}
    &
    W_Q' = \left(\begin{matrix}
        W_Q                         & 0^{\Floating^{d \times 1}} & 0^{\Floating^{d \times 1}} & 0^{\Floating^{d \times 1}} \\
        0^{\Floating^{d \times 1}} & 0 & 0 & 0 \\
        0^{\Floating^{d \times 1}} & 1 & 0 & 0 \\
        0^{\Floating^{d \times 1}} & 1 & 0 & 0 \\
    \end{matrix}\right)
    &
    W_K' = \left(\begin{matrix}
        W_K                         & 0^{\Floating^{d \times 1}} & 0^{\Floating^{d \times 1}}       & 0^{\Floating^{d \times 1}} \\
        0^{\Floating^{d \times 1}} & 0 & 0       & 0 \\
        0^{\Floating^{d \times 1}} & 0 & \minrep & 0 \\
        0^{\Floating^{d \times 1}} & 0 & 0       & \minrep \\
    \end{matrix}\right)
\end{align*}
%
Then, 
%
\begin{align}
    \langle q'_i, k'_j \rangle &= \sum_{x=1}^{d+3} (q'_l)_x (k'_i)_x = \langle q_i, k_j \rangle + (q'_l)_{d+1} (k'_i)_{d+1} + (q'_l)_{d+2} (k'_i)_{d+2} + (q'_l)_{d+3} (k'_i)_{d+3} \nonumber \\
    &
    = \langle q_i, k_j \rangle + 0 + (q'_l)_{d+2} (k'_i)_{d+2} + (q'_l)_{d+3} (k'_i)_{d+3} \nonumber \\
    &
    = \langle q_i, k_j \rangle + (q'_l)_{d+2} (k'_i)_{d+2} + (q'_l)_{d+3} (k'_i)_{d+3} \nonumber \\
    & 
    = \begin{cases}
        \langle q_i, k_j \rangle + 1 \cdot \minrep + 1 \cdot \minrep & \text{ if } i = l    \land j \in M \\
        \langle q_i, k_j \rangle + 1 \cdot 0       + 1 \cdot 0       & \text{ if } i = l    \land j \notin M \\
        \langle q_i, k_j \rangle + 0 \cdot \minrep + 0 \cdot \minrep & \text{ if } i \neq l \land j \in M \\
        \langle q_i, k_j \rangle + 0 \cdot 0       + 0 \cdot 0       & \text{ if } i \neq l \land j \notin M \\
    \end{cases} \nonumber \\
    & 
    = \begin{cases}
        \langle q_i, k_j \rangle + \minrep + \minrep & \text{ if } i = l    \land j \in M \\
        \langle q_i, k_j \rangle + 0       + 0       & \text{ if } i = l    \land j \notin M \\
        \langle q_i, k_j \rangle + 0       + 0       & \text{ if } i \neq l \land j \in M \\
        \langle q_i, k_j \rangle + 0       + 0       & \text{ if } i \neq l \land j \notin M \\
    \end{cases}\nonumber \\
    & 
    = \begin{cases}
        \minrep                  & \text{ if } i = l    \land j \in M \\
        \langle q_i, k_j \rangle & \text{ otherwise } 
    \end{cases} \nonumber
\end{align}
% 
Namely, $\langle q'_i, k'_j \rangle = \minrep$ if $i = l$ and $j \in M$, but $\langle q'_i, k'_j \rangle = \langle q_i, k_j \rangle$ otherwise.



\section{Linear transformations}\label{appendix:linear-transf}

We show a mechanism to replace a flagged vector with the result of applying the linear transformation $M$ to it.
Let $f$ denote the index of the flagged vector.
We perform the operation with two consecutive $\ATTN$ layers, in which $h_f$ only attends to itself.
Therefore, after these two layers $h_f$ will be replaced with $Mh_f$.

We need to double the embedding size from $d$ to $2d$ (plus some flags) since we need to store the result of the transformation in some coordinates (the added $d+1$ to $2d$) as a middle step.
In the first $\ATTN$ layer, we compute the transformation saving its result in the coordinates $d+1$ to $2d$.
Then, in the second layer we move the result to the right place.

\[
    \widetilde{h_f} = \left( \begin{matrix} (h_f)_1    \\ \vdots \\ (h_f)_d    \\ 0      \\ \vdots \\ 0       \end{matrix} \right) 
    \xrightarrow{\text{first $\ATTN$ layer}}
    \widetilde{h_f} = \left( \begin{matrix} 0      \\ \vdots \\ 0      \\ (Mh_f)_1 \\ \vdots \\ (Mh_f)_d \end{matrix} \right)
    \xrightarrow{\text{second $\ATTN$ layer}}
    \widetilde{h_f} = \left( \begin{matrix} (Mh_f)_1 \\ \vdots \\ (Mh_f)_d \\ 0      \\ \vdots \\ 0      \end{matrix} \right)
\]

We assume that the $d+1$ to $2d$ coordinates (used as memory) are set to $0$ by the layers preceding this mechanism
This can be achieved maintaining them as $0$ since the beginning or transformed to $0$ with the primitive described in Section \ref{const-fun}.

\medskip

Let us construct each of the $\ATTN$ layers.
As said before, in both layers we will need for $\widetilde{h_f}$ to only attend to itself.
Therefore, we define $W_Q$ and $W_K$ using the strategy of Section \ref{ignore-vecs}.

Recall from the definition of the $\ATTN$ mechanism (\ref{alg:def_attn}) that the output of the $\ATTN$ layer is $h'_i = W_O \sum_{j=1}^a (s_i)_j v_j$. 
By this point, $h'_f$ is looking like \[ h'_f = W_O \sum_{j=1}^a (s_f)_j v_j = W_O v_f = W_O (W_V \widetilde{h_f}) \]

\begin{itemize}
\item For the first layer, defining $W_O$ and $W_V$ as follows
%
\begin{align*}
    W_V &= id
    & 
    W_O &= \left(\begin{matrix}
                -id^{\Rdd}                 &0^{\Rdd}                   \\
                M                          &0^{\Rdd}                   
          \end{matrix}\right)
\end{align*}
%
makes us get $ h'_f = (-(h_f)_1, \dots, -(h_f)_d, (Mh_f)_1, \dots, (Mh_f)_d)$, which added to $\widetilde{h_f}$ gives us what we were looking for after the first layer.


\item For the second $\ATTN$ layer, we use the same mechanism for making $\widetilde{h_f}$ only attend to itself.
And defining $W_V$ and $W_O$ as follows
%
\begin{align*}
    W_V &= id
    & 
    W_O &= \left(\begin{matrix}
                0^{\Rdd}                   &id^{\Rdd}                  \\
                0^{\Rdd}                   &-id^{\Rdd}                 \\
          \end{matrix}\right)
\end{align*}
%
we get $ h'_f = ((Mh_f)_1, \dots, -(Mh_f)_d, -(Mh_f)_1, \dots, -(Mh_f)_d)$, which added to $\widetilde{h_f}$ gives us what we were looking for after the second layer.

\end{itemize}

\bigskip

Note that, after these two layers, the values of $h_i$ for all $i \neq f$ have changed in ways that have not been studied in this paper, but this is not relevant to the way we use the primitive.





\section{Implementation of the special operation}\label{appendix:special-op}

This operation is not a linear transformation; therefore, we need to define a new technique.
Observe that one of $\mathbb{I}_a$ and $\mathbb{I}_b$ equals $1$ and the other equals $0$.
Then, $\mathbb{I}_b = 1 - \mathbb{I}_a$ and $\mathbb{I}_a = 1 - \mathbb{I}_b$.

\begin{equation*}
    \left(\begin{matrix}
        a \\ b \\ \mathbb{I}_a \\ \mathbb{I}_b
    \end{matrix}\right)
    \xrightarrow[\text{operation}]{\text{Special}}
    \left(\begin{matrix}
        a \cdot \mathbb{I}_a \\ b \cdot \mathbb{I}_b\\ 0 \\ 0
    \end{matrix}\right)
\end{equation*}

This operation is performed with three layers of $\FF$.
Naturally, after each layer a new vector is added to the original.
The vectors we add are the following:

\begin{equation*}
    \left(\begin{matrix}
        a \\ b \\ \mathbb{I}_a \\ \mathbb{I}_b 
    \end{matrix}\right)
    \xrightarrow[\text{operation}]{\text{Special}}
    \left(\begin{matrix}
        a            &+& \maxrep \cdot \mathbb{I}_b &+& \maxrep \cdot \mathbb{I}_b &+& \minrep \cdot \mathbb{I}_b \\
        b            &+& \maxrep \cdot \mathbb{I}_a &+& \maxrep \cdot \mathbb{I}_a &+& \minrep \cdot \mathbb{I}_a \\ 
        \mathbb{I}_a &+& 0                          &+& 0                          &+& -\mathbb{I}_a\\
        \mathbb{I}_b &+& 0                          &+& 0                          &+& -\mathbb{I}_b\\
    \end{matrix}\right)    
\end{equation*}


Note that if $\mathbb{I}_a = 1$, then $\mathbb{I}_b = 0$, so the first coordinate isn't modified and the second is changed to $0$ because of the finite representation properties.
In the same way, if $\mathbb{I}_a = 0$, then $\mathbb{I}_b = 1$, so the second coordinate isn't modified and the first is changed to $0$.

\medskip

Let us show how to instantiate the $\FF$ layers.
The first two have parameters:

\begin{align*}
    & W_1 = id
    & W_2  = \left(\begin{matrix}
        0  & 0  & 0        & \maxrep   \\
        0  & 0  & \maxrep  & 0         \\
        0  & 0  & 0        & 0         \\
        0  & 0  & 0        & 0         \\
    \end{matrix}\right) &
    & b_1 = b_2 = \left(\begin{matrix}
        0 \\
        0 \\
        0 \\
        0
    \end{matrix}\right) 
\end{align*}


\noindent Notice that the output of these $\FF$ layer when the input is $(x, y, \mathbb{I}_a, \mathbb{I}_b)$ will be $(\maxrep \cdot \mathbb{I}_b, \maxrep \cdot \mathbb{I}_a, 0, 0)$.

The third $\FF$ layer has parameters:

\begin{align*}
    & W_1 = id
    & W_2  = \left(\begin{matrix}
        0  & 0  & 0        & \minrep   \\
        0  & 0  & \minrep  & 0         \\
        0  & 0  & -1       & 0         \\
        0  & 0  & 0        & -1        \\
    \end{matrix}\right) &
    & b_1 = b_2 = \left(\begin{matrix}
        0 \\
        0 \\
        0 \\
        0
    \end{matrix}\right) 
\end{align*}

\noindent which, when the input is $(x, y, \mathbb{I}_a, \mathbb{I}_b)$, outputs $(\minrep \cdot \mathbb{I}_b, \minrep \cdot \mathbb{I}_a, -\mathbb{I}_a, -\mathbb{I}_b)$




\bigrafa{intenté escribir cómo quedaban las $\FF$ instanciadas pero queda medio ilegible. No sé si vale la pena}




\section{Repeat tokens}\label{appendix:repeat-token}

We show how to make a transformer repeat a token until we decide it to make it stop.
This construction is needed to argue that for every transformer $\TF$ there is a transformer $\TF'$ of the same $\CoTtndn$ class that computes the same language and consumes all of its budget.

Assume that $\TF$ is normalized and let us define $\TF'$.
$\TF'$ is identical to $\TF$ except for the following modifications.
Let $\sigma$ be the token to repeat.
We flag the vectors coming from $\sigma$ with the token embedding.
Since we are replicating a token that was printed, i.e., it is not part of the input, its corresponding vector is the last one.
This is because we begin repeating it after the first time it is printed.


Since we flag the vectors coming from the token $\sigma$ we can apply a constant function only to them.
This function is applied in the last layer of the body of $\TF'$.
Since we assume $\TF$ was normalized, the $\OUTPUT$ matrix is just a projection.
Then, the value of the constant function is what is outputted as the probability distribution (after $\softmax$).
If the constant function has value $\minrep$ in every coordinate except in the one associated with the probability of $\sigma$, where it has $\maxrep$, the distribution outputted by the distribution generator will have all the probability concentrated in $\sigma$.
Thus, $\TF'$ prints $\sigma$.

\medskip

Observe that this only works if the token does not appear in the input word.
This case can be handled by exploiting the non-uniformity of the model.
If the input word is of size $n$, with the $\PE$ we can flag the first $n$ vectors.
Then, we can use that flag to set to $0$ the coordinate used by the $\TE$ for flagging the token to be repeated, i.e., turning of the flag.
We make them $0$ using a constant function as described in the Section \ref{const-fun}.


